% ConjugAI — diagramas para inclusão em trabalho académico
% Compilar: pdflatex diagramas-conjugai.tex
% Requer: distribuição TeX com TikZ (texlive-full, MiKTeX, etc.)

\documentclass[a4paper,11pt]{article}
\usepackage[utf8]{inputenc}
\usepackage[T1]{fontenc}
\usepackage[portuguese]{babel}
\usepackage{microtype}
\usepackage[hyphens]{url}
\makeatletter
\g@addto@macro\UrlBreaks{\do\/\do\_}
\makeatother
\usepackage{tikz}
\usetikzlibrary{positioning,shapes.geometric,arrows.meta,fit,backgrounds,calc}

\tikzset{
  proc/.style={rectangle, draw, rounded corners, minimum width=2.8cm, align=center, inner sep=3pt},
  dec/.style={diamond, draw, aspect=2.2, align=center, inner sep=1pt, font=\small},
  err/.style={rectangle, draw, rounded corners, fill=red!8, align=center, inner sep=3pt, font=\small},
  lbl/.style={font=\footnotesize, midway, fill=white, inner sep=1pt},
  >=Stealth
}

\title{ConjugAI --- Figuras (fluxos e arquitetura)}
\author{}
\date{}

\begin{document}
\maketitle

\noindent
Estas figuras correspondem aos diagramas Mermaid do repositório
(\url{diagrama.html}, \url{docs/diagrama.md}).
Pode copiar cada ambiente \texttt{tikzpicture} para o seu documento ou incluir este PDF com \verb|\includegraphics|.

\section{Pipeline \texttt{analisarFrase}}

\begin{figure}[htbp]
  \centering
  \resizebox{0.95\textwidth}{!}{%
  \begin{tikzpicture}[node distance=6mm, auto]
    \node[proc] (a) {Texto bruto};
    \node[proc, below=of a] (b) {\texttt{tokenize}};
    \node[dec, below=of b] (c) {Tokens vazios?};
    \node[err, below left=8mm and 4mm of c, align=center, font=\scriptsize] (err) {Erro:\\Digite ou selecione uma frase};
    \node[proc, below right=8mm and 4mm of c] (d) {\texttt{detectarSujeito}};
    \node[proc, below=of d] (e) {\texttt{detectarTempo}};
    \node[proc, below=of e] (f) {\texttt{extrairVerbo}};
    \node[dec, below=of f] (g) {Lema verbal?};
    \node[err, below left=8mm and 4mm of g] (err2) {Erro: verbo não identificado};
    \node[proc, below right=8mm and 4mm of g] (h) {\texttt{conjugarTempo}};
    \node[dec, below=of h] (i) {Forma no tempo?};
    \node[err, below left=8mm and 4mm of i] (err3) {Erro: não foi possível conjugar};
    \node[proc, below right=8mm and 4mm of i] (j) {\texttt{corrigir}};
    \node[proc, below=of j] (k) {Frase corrigida};

    \draw[->] (a) -- (b);
    \draw[->] (b) -- (c);
    \draw[->] (c) -| node[near start, left, font=\scriptsize] {sim} (err);
    \draw[->] (c) -| node[near start, right, font=\scriptsize] {não} (d);
    \draw[->] (d) -- (e);
    \draw[->] (e) -- (f);
    \draw[->] (f) -- (g);
    \draw[->] (g) -| node[near start, left, font=\scriptsize] {não} (err2);
    \draw[->] (g) -| node[near start, right, font=\scriptsize] {sim} (h);
    \draw[->] (h) -- (i);
    \draw[->] (i) -| node[near start, left, font=\scriptsize] {não} (err3);
    \draw[->] (i) -| node[near start, right, font=\scriptsize] {sim} (j);
    \draw[->] (j) -- (k);
  \end{tikzpicture}%
  }
  \caption{Fluxo principal desde o texto bruto até \texttt{correcao} (\texttt{vendors/conjugai-core/index.ts}).}
  \label{fig:pipeline-analisar}
\end{figure}

\section{Função \texttt{extrairVerbo}}

\begin{figure}[htbp]
  \centering
  \resizebox{0.92\textwidth}{!}{%
  \begin{tikzpicture}[node distance=7mm, auto]
    \node[proc] (x) {\texttt{extrairVerbo(tokens)}};
    \node[dec, below=of x] (a) {há infinitivo\\\texttt{-ar/-er/-ir/-pôr}?};
    \node[proc, below left=10mm and 2mm of a] (b) {retorna 1.\textsuperscript{o} infinitivo};
    \node[dec, below right=10mm and 2mm of a] (c) {1.\textsuperscript{o} token é\\\texttt{vou/vais/vai/...}\\e há infinitivo?};
    \node[proc, below left=10mm and 2mm of c] (d) {retorna \texttt{ir}};
    \node[proc, below right=10mm and 2mm of c] (e) {\texttt{detectarVerboPorDicionario}\\(filtro funcional)};
    \node[dec, below=of e] (f) {lema?};
    \node[proc, below left=10mm and 2mm of f] (g) {retorna infinitivo};
    \node[proc, below right=10mm and 2mm of f] (h) {\texttt{null}};

    \draw[->] (x) -- (a);
    \draw[->] (a) -| node[near start, left, font=\scriptsize] {sim} (b);
    \draw[->] (a) -| node[near start, right, font=\scriptsize] {não} (c);
    \draw[->] (c) -| node[near start, left, font=\scriptsize] {sim} (d);
    \draw[->] (c) -| node[near start, right, font=\scriptsize] {não} (e);
    \draw[->] (e) -- (f);
    \draw[->] (f) -| node[near start, left, font=\scriptsize] {sim} (g);
    \draw[->] (f) -| node[near start, right, font=\scriptsize] {não} (h);
  \end{tikzpicture}%
  }
  \caption{Identificação do lema verbal (\texttt{conjugador.ts}).}
  \label{fig:extrair-verbo}
\end{figure}

\section{Sujeito e tempo verbal}

\begin{figure}[htbp]
  \centering
  \resizebox{0.98\textwidth}{!}{%
  \begin{tikzpicture}[node distance=7mm]
    \node[dec] (sc) {Prefixo \emph{X e Y} antes do verbo?};
    \node[proc, below left=14mm and 6mm of sc, align=center, font=\small] (sc1) {Nós / Vocês / Eles\\(\texttt{detectarSujeitoComposto})};
    \node[dec, below right=14mm and 6mm of sc] (s0) {Eu + mamãe ou papai?};
    \node[proc, below left=12mm and 4mm of s0] (sn) {Nós --- 1.\textsuperscript{a} plural};
    \node[proc, below right=12mm and 4mm of s0, align=center, font=\small] (s1) {Ordem: nós, eles, ela, ele, eu, padrão eu};
    \draw[->] (sc) -| node[near start, left, font=\scriptsize] {sim} (sc1);
    \draw[->] (sc) -| node[near start, right, font=\scriptsize] {não} (s0);
    \draw[->] (s0) -| node[near start, left, font=\scriptsize] {sim} (sn);
    \draw[->] (s0) -| node[near start, right, font=\scriptsize] {não} (s1);
  \end{tikzpicture}%
  }
  \caption{Ordem em \texttt{sujeito.ts}: composto antes das regras simples.}
  \label{fig:sujeito}
\end{figure}

\begin{figure}[htbp]
  \centering
  \resizebox{0.98\textwidth}{!}{%
  \begin{tikzpicture}[node distance=8mm]
    \node[dec] (ovr) {Tempo explícito\\\texttt{tempo:x} / \texttt{[tempo=x]}?};
    \node[dec, below left=12mm and 18mm of ovr] (j1) {Tem \emph{ontem}+\emph{já}?};
    \node[dec, below right=12mm and 18mm of ovr] (j2) {Tem \emph{amanhã}+\emph{já}?};
    \node[proc, below=11mm of j1, font=\small] (ppc) {Pretérito perfeito composto};
    \node[proc, below=11mm of j2, font=\small] (fc) {Futuro composto};
    \node[dec, below=16mm of ovr] (s1) {Tem \emph{se}/\emph{quando}/\emph{talvez}/\emph{que}?};
    \node[proc, below=11mm of s1, font=\small] (sub) {Subjuntivo (heurística)};
    \node[dec, below=11mm of sub] (t0) {Tem \emph{amanhã}?};
    \node[dec, below left=12mm and 17mm of t0] (t1) {Tem \emph{ontem}?};
    \node[dec, below right=12mm and 17mm of t0] (p1) {1.\textsuperscript{o} token vou/vai…?};
    \node[proc, below=11mm of t1, font=\small] (tp) {Pretérito perfeito};
    \node[proc, below=11mm of p1, font=\small] (tf) {Futuro do presente};
    \node[proc, below=27mm of t0, font=\small] (tn) {Presente do indicativo};
    \node[proc, right=20mm of ovr, font=\small] (texp) {Usa tempo explícito};

    \draw[->] (ovr.east) -- node[midway, above, font=\scriptsize] {sim} (texp.west);
    \draw[->] (ovr.south west) -- node[midway, left, font=\scriptsize] {não} (j1.north);
    \draw[->] (ovr.south east) -- node[midway, right, font=\scriptsize] {não} (j2.north);
    \draw[->] (j1.south) -- node[midway, left, font=\scriptsize] {sim} (ppc.north);
    \draw[->] (j2.south) -- node[midway, right, font=\scriptsize] {sim} (fc.north);
    \draw[->] (j1.east) -- ++(6mm,0) |- node[pos=0.35, above, font=\scriptsize] {não} (s1.west);
    \draw[->] (j2.west) -- ++(-6mm,0) |- node[pos=0.35, above, font=\scriptsize] {não} (s1.east);
    \draw[->] (s1.south) -- node[midway, left, font=\scriptsize] {sim} (sub.north);
    \draw[->] (s1.south) -- ++(0,-4mm) node[right, font=\scriptsize] {não} -- ++(0,-2mm) -- (t0.north);
    \draw[->] (sub.south) -- (t0.north);
    \draw[->] (t0.south west) -- node[midway, left, font=\scriptsize] {não} (t1.north);
    \draw[->] (t0.south east) -- node[midway, right, font=\scriptsize] {sim} (p1.north);
    \draw[->] (t1.south) -- node[midway, left, font=\scriptsize] {sim} (tp.north);
    \draw[->] (p1.south) -- node[midway, right, font=\scriptsize] {não} (tf.north);
    \draw[->] (t1.east) -- ++(6mm,0) |- node[pos=0.30, above, font=\scriptsize] {não} (tn.west);
    \draw[->] (p1.west) -- ++(-6mm,0) |- node[pos=0.30, above, font=\scriptsize] {sim} (tn.east);
  \end{tikzpicture}%
  }
  \caption{Seleção de tempo em \texttt{tempo.ts}: override explícito e heurísticas por marcadores.}
  \label{fig:tempo}
\end{figure}

\section{Arquitetura da aplicação (demo CAA)}

\begin{figure}[htbp]
  \centering
  \resizebox{\textwidth}{!}{%
  \footnotesize
  \begin{tikzpicture}[
    node distance=4mm,
    box/.style={rectangle, draw, rounded corners, align=center, inner sep=2pt, font=\scriptsize},
    grp/.style={rectangle, draw, dashed, inner sep=6pt}
  ]
    \node[box] (aside) {Lista / select\\de exemplos};
    \node[box, right=of aside] (in) {Textarea\\entrada};
    \node[box, right=of in] (btn) {Analisar / Limpar};
    \node[box, below=8mm of in] (steps) {4 passos:\\tokens $\rightarrow$ sujeito $\rightarrow$ tempo $\rightarrow$ regra};
    \node[box, right=of steps] (out) {Saída\\destacada};

    \node[box, below=14mm of steps] (af) {\texttt{analisarFrase}};
    \node[box, left=6mm of af] (tok) {\texttt{tokenize}};
    \node[box, right=6mm of af] (sub) {\texttt{detectarSujeito}};
    \node[box, right=of sub] (tmp) {\texttt{detectarTempo}};
    \node[box, right=of tmp] (cj) {\texttt{conjugar} + \texttt{corrigir}};

    \node[box, below=10mm of af] (an) {\texttt{analyze}};
    \node[box, right=of an] (anim) {\texttt{runStepAnimation}};

    \begin{scope}[on background layer]
      \node[grp, fit=(aside)(in)(btn)(steps)(out), label=above:\textbf{Interface (demo CAA)}] {};
      \node[grp, fit=(tok)(af)(sub)(tmp)(cj), label=above:\textbf{conjugai-core.js}] {};
      \node[grp, fit=(an)(anim), label=below:\textbf{app.js}] {};
    \end{scope}

    \draw[->] (aside) -- (in);
    \draw[->] (in) -- (btn);
    \draw[->] (btn.south) -- ++(0,-4mm) -| (an.north);
    \draw[->] (an) -- (af);
    \draw[->] (af) -- (tok);
    \draw[->] (tok) -- (sub);
    \draw[->] (sub) -- (tmp);
    \draw[->] (tmp) -- (cj);
    \draw[->] (an) -- (anim);
    \draw[->] (anim.north) |- (steps.west);
    \draw[->] (anim.north) -| (out.south);
  \end{tikzpicture}%
  }
  \caption{Componentes da SPA e chamadas ao núcleo (simplificado).}
  \label{fig:arquitetura}
\end{figure}

\section{Sequência --- demonstração ao vivo}

\begin{figure}[htbp]
  \centering
  \resizebox{0.98\textwidth}{!}{%
  \begin{tikzpicture}[
    >=Stealth,
    actor/.style={rectangle, draw, rounded corners, align=center, font=\scriptsize, minimum width=3.1cm, minimum height=7mm},
    msg/.style={->, thick},
    rmsg/.style={->, thick, dashed}
  ]
    % Colunas (mesma linha de base)
    \node[actor] (u) at (0,0) {Utilizador};
    \node[actor] (s) at (4.2,0) {Select / Textarea};
    \node[actor] (ui) at (8.7,0) {Painel de passos + saída};
    \node[actor] (core) at (13.6,0) {\texttt{ConjugaiCore.analisarFrase}};

    % Lifelines (tempo desce no eixo $y$)
    \draw[gray!60, dashed] (u.south) -- ++(0,-6.2);
    \draw[gray!60, dashed] (s.south) -- ++(0,-6.2);
    \draw[gray!60, dashed] (ui.south) -- ++(0,-6.2);
    \draw[gray!60, dashed] (core.south) -- ++(0,-6.2);

    % Ancoragem explícita com calc ($\cdot$): evita setas diagonais espúrias
    \def\ya{-0.55}
    \def\yb{-1.35}
    \def\yc{-2.15}
    \def\yd{-2.95}
    \def\ye{-3.75}
    \def\yf{-4.85}

    % 1: utilizador edita o select/textarea
    \draw[msg]
      ($(u.south)+(0,\ya)$) -- ($(s.south)+(0,\ya)$)
      node[midway, above, font=\tiny, inner sep=1pt] {escolhe exemplo / edita texto};

    % 2: utilizador dispara análise na UI
    \draw[msg]
      ($(u.south)+(0,\yb)$) -- ($(ui.south)+(0,\yb)$)
      node[midway, above, font=\tiny, inner sep=1pt] {clica Analisar};

    % 3: UI envia texto ao núcleo
    \draw[msg]
      ($(ui.south)+(0,\yc)$) -- ($(core.south)+(0,\yc)$)
      node[midway, above, font=\tiny, inner sep=1pt] {texto bruto};

    % 4: retorno do núcleo para a UI
    \draw[rmsg]
      ($(core.south)+(0,\yd)$) -- ($(ui.south)+(0,\yd)$)
      node[midway, above, font=\tiny, inner sep=1pt] {\texttt{ResultadoAnalise}};

    % Nota: processamento local na UI (sem seta cruzada)
    \node[anchor=west, font=\tiny, align=left] at ($(ui.south)+(0.35,\ye)$)
      {$\hookrightarrow$ animação 4 etapas ($\sim$1{,}8 s)};

    % 5: resultado apresentado ao utilizador
    \draw[msg]
      ($(ui.south)+(0,\yf)$) -- ($(u.south)+(0,\yf)$)
      node[midway, above, font=\tiny, inner sep=1pt] {frase corrigida (saída)};
  \end{tikzpicture}
  }
  \caption{Troca de mensagens entre utilizador, entrada, UI e núcleo (conceitual).}
  \label{fig:sequencia}
\end{figure}

\end{document}
